\begin{afterword}
\summaryhead

The author reveals that the early ``mysterious answer'' mentioned in earlier posts concerned morality as well as consciousness. In 1997 the young author, Eliezer1997, argued that if life has no meaning every outcome has utility zero, so that case cancels out of any expected-utility calculation; one can act as if life is meaningful, and since humans cannot find the meaning, build a superintelligence to find it. The objection that no ``ought'' follows from logic was met by grouping three ``hard problems,'' consciousness, existence and morality, as too hard for humans and perhaps needing new physics.

The author diagnoses the errors: utility treated as an absolute quantity, ``should'' treated as whatever compels any mind and equated with ``right,'' and the supposed proof used as a reason not to define ``intelligence'' or ``meaning.'' The young author was skilled at refuting others but did not turn that skill on this position. The lessons: no clever justification for relaxing rigor protects you, since Nature grades only results; vague concepts may guide thought but must not be built on; and confusion cannot be worked around, only dissolved.

\responsehead

The post is a confession, and on the points I could check it is an accurate one. The author declines to reread the old writings and warns that memory is reconstructed. Yet the argument given here matches the author's ``Frequently Asked Questions about the Meaning of Life'' of 1999. There the meaningless case is every goal having zero desirability; it drops out; one may ``assume that life has meaning and work from there''; and the Singularity is ``the interim meaning of life.'' The same text calls its conclusion ``an informed guess,'' the kind of ``best guess'' that, the post now says, excused nothing.

The diagnosis is also sound where it is technical. Utilities are fixed only up to an affine transformation, so a case in which every outcome is equally valuable drops out whatever that value is. A recent book on moral uncertainty, by MacAskill, Bykvist and Ord, treats such a uniform theory the same way, and on that treatment the cancelling step itself is valid; the post's charge is that it was applied to a ``meaning'' left undefined. The deeper error it names, taking what would compel any mind to be what ``right'' means, is the one the book argued against in ``No Universally Compelling Arguments,'' and the notes there place it among standard views.

The first lesson claims more than the case can show. ``No matter how clever the justification for relaxing your standards \ldots\ it will blow your foot off just the same'' is a rule for every case, drawn from one case, which shows that relaxing rigor can go wrong, not that it always does. The second rule ends narrower than it begins. The flat ``No, you don't use the best concepts you can use at the time'' becomes, a few paragraphs later: vague concepts may guide the search, but one does not build on them.

The reply to the young author's reason for haste rests on a premise the post asserts: that a best guess fails as surely as no attempt (``Either you match it, or you fail''). The young author's reason was a cost of waiting, nanoweapons arriving first. Nature's indifference settles that only if the premise holds. The book argues for it in ``Value is Fragile,'' written four months later.

In short: an accurate account of the author's 1997 argument, as the author's own 1999 text confirms, with a sound diagnosis; its lessons are stated for every case on the strength of one, and its answer to the case for haste rests on a premise the post asserts and the book argues later.
\end{afterword}
